% ============================================================
% 中文摘要
% 格式要求（2025版）：
%   标题：黑体三号，居中
%   正文：小四号宋体，两端对齐，300-600汉字
%   关键词：不超过5个，关键词之间用分号（；）间隔
% ============================================================
\cleardoublepage{}
\markboth{摘要}{}
\begin{center}
    \heiti\zihao{3} 摘\hspace{1em}要
\end{center}

% 【请将以下内容替换为你自己的中文摘要正文】
随着……的发展，……问题日益突出。本文提出了一种新型……方法，融合了……关键技术，通过……实验设计，在……数据集上验证了方法的有效性，取得了……指标的性能提升。研究表明，该方法能够……，为……问题提供了新的解决方案。

% 关键词单独成段，标题加粗，词间用分号
\vspace{\baselineskip}
\noindent{\heiti 关键词：}关键词1；关键词2；关键词3；关键词4；关键词5

% ============================================================
% 英文摘要
% 格式要求（2025版）：
%   标题：Arial体三号，居中（模板中 \sffamily 对应 Arial）
%   正文：小四号 Times New Roman，两端对齐
%   Keywords：与中文关键词对应，词间用分号（;）间隔
%   内容须与中文摘要对应
% ============================================================
\cleardoublepage{}
\markboth{Abstract}{}
\begin{center}
    \sffamily\zihao{3} Abstract
\end{center}

% 【请将以下内容替换为你自己的英文摘要正文】
With the development of ..., the problem of ... has become increasingly prominent.
This paper proposes a novel ... method that integrates ... key techniques.
Through ... experimental design, its effectiveness is validated on ... datasets,
achieving performance improvements in ... metrics.
The results demonstrate that this approach can ..., providing a new solution to ...

% Keywords 单独成段
\vspace{\baselineskip}
\noindent\textbf{Keywords}: Keyword1; Keyword2; Keyword3; Keyword4; Keyword5
